\section{Preguntas de análisis}

\subsection{Exclusión por página y operaciones sobre varias páginas}

La exclusión por página garantiza que, mientras un nodo tiene permiso de
escritura sobre una página, ningún otro la lee ni la escribe, y que lo que
publique será visible para quien adquiera después con una versión mayor. Es
una garantía \emph{por página y por intervalo}: vale entre el
\code{DSM\_acquire} y el \code{DSM\_release} de esa página.

No garantiza nada sobre una operación que abarque varias páginas. Si un nodo
adquiere la página 0, la modifica, la libera, y después adquiere la página 1,
otro nodo pudo haber escrito la página 1 en el medio —o haber leído la 0 ya
modificada y la 1 todavía vieja. Las dos escrituras son atómicas por separado
y el conjunto no lo es. El protocolo no ofrece ninguna forma de adquirir dos
páginas como una unidad, y adquirirlas de a una en orden fijo evita el
interbloqueo pero no produce atomicidad: solo evita que dos nodos se esperen
en círculo.

Una barrera final no corrige incrementos perdidos porque una barrera es una
afirmación sobre el \emph{tiempo}, no sobre el \emph{contenido}. Dice que
todos llegaron a este punto; no dice qué hizo cada uno antes de llegar. Si dos
nodos leyeron el valor 7, le sumaron uno y publicaron 8, el 8 es el valor
confirmado y la barrera posterior lo deja exactamente así. El incremento se
perdió en el momento en que el segundo nodo leyó antes de que el primero
publicara, y nada que ocurra después puede saber que ese 9 debió existir.

La forma de no perderlo es que la lectura y la escritura ocurran dentro del
mismo intervalo de escritura: adquirir en modo \code{WRITE}, leer, sumar y
liberar. Entonces el segundo nodo no puede leer hasta que el primero publicó,
y el protocolo lo impide por construcción, no por disciplina del programador.

\subsection{El coordinador central}

Para razonar sobre corrección, el coordinador central tiene una ventaja que
cuesta sobrevalorar: \emph{existe un único orden}. Todas las decisiones se
toman en un solo hilo, una por vez, sobre un estado que nadie más toca. Eso
convierte afirmaciones que serían difíciles —``ningún lector observó la página
mientras el escritor la tenía''— en propiedades de una máquina de estados que
se pueden comprobar después de cada transición y que se pueden leer en una
secuencia de filas numeradas.

También hace que la verificación sea posible desde afuera. El registro del
coordinador es la historia completa del sistema: no hay eventos que ocurrieran
en otro lado y que haya que correlacionar por reloj.

El costo está en el mismo lugar que la ventaja. El coordinador es un cuello de
botella de dos maneras. Por trabajo, todas las solicitudes pasan por un hilo,
aunque en este sistema ese hilo solo compara máscaras de bits y el límite real
está en la red. Y por tráfico: cada página publicada viaja al coordinador y
vuelve a salir hacia quien la pida, de modo que una transferencia entre dos
nodos cuesta dos veces la página en vez de una. Un protocolo con directorio
distribuido mandaría la página del dueño anterior al nuevo directamente.

Como punto único de fallo es total y está asumido: si el coordinador muere, la
sesión muere. No hay réplica del directorio ni de los datos confirmados, de
modo que no hay ningún estado del que otro proceso pudiera continuar. Los
trabajadores también son puntos únicos en un sentido más débil: la caída de
uno aborta la sesión, porque las barreras esperan a todos y una página que él
tuviera en escritura no tiene dueño que la publique.

\subsection{Tráfico cuando datos independientes comparten página}

El experimento de falso compartimiento responde esto con números. Con dos
nodos que incrementan contadores que nadie más mira, ponerlos en la misma
página sube el tráfico de coherencia de 433\,280 a 832\,640 bytes y las
páginas transferidas de 102 a 198; las invalidaciones pasan de cero a 96. Con
cuatro nodos, de cero a 198,6 invalidaciones y de 866\,560 a
1\,692\,736 bytes.
El trabajo de la aplicación es idéntico y los bytes de coordinación también:
lo único que cambió es un desplazamiento dentro de la memoria.

La razón es que el protocolo no sabe qué parte de la página le importa a
quién. Cuando un nodo pide escribir, se invalida la página completa en todos
los que tengan copia, aunque la zona que van a leer no se haya tocado. La
unidad de coherencia es la página y la granularidad del dato es irrelevante
para ella.

Replicar lecturas conviene cuando la página se lee mucho más de lo que se
escribe. Varios lectores coexisten sin costo entre ellos: cada uno paga la
transferencia una vez y después lee a velocidad de memoria mientras conserve
el permiso. El experimento de localidad lo muestra en la dirección contraria:
el costo no está en leer sino en volver a pedir permiso.

Deja de convenir en cuanto aparece una escritura, y el costo que agrega no es
el de la escritura sino el de las copias. Escribir una página que nadie más
tiene es un mensaje y una transferencia. Escribir una que tienen otros tres
cuesta además tres invalidaciones, tres confirmaciones, y la próxima lectura
de cada uno vuelve a traer los 4096 bytes. El costo de escribir crece con la
cantidad de copias, no con el tamaño del dato que se escribió.

\subsection{Los costos que la interfaz oculta}

La interfaz de memoria compartida oculta que cada acceso puede ser una
conversación. \code{DSM\_acquire} parece una operación local y puede costar
una solicitud, una ronda de invalidaciones con sus confirmaciones, una espera
en cola detrás de otros nodos y una transferencia de 4096 bytes. El puntero
que devuelve se usa con instrucciones comunes, y esa es exactamente la
apariencia que esconde el costo.

También oculta el momento. El dato que se lee es el de la última publicación,
que puede ser de hace un microsegundo o de hace un minuto; nada en la interfaz
distingue esos dos casos. Y oculta la granularidad: el programador declara
variables, el protocolo mueve páginas, y el falso compartimiento vive
enteramente en la distancia entre esas dos unidades.

Nuestro contrato mitiga esto no por ser más barato, sino por ser explícito. El
costo sigue estando, pero aparece donde el programador escribió una llamada.
Eso es lo que permite medirlo por intervalo y atribuirlo a una línea de
código.

Recuperar la transparencia por fallos de página —que la aplicación use
punteros comunes sin pedir permiso— requeriría bastante más que interceptar la
señal. Haría falta delimitar la región de memoria que es DSM y protegerla con
\code{mprotect}, para que el acceso a una página sin permiso produzca un
fallo. Habría que distinguir ese fallo de un \code{SIGSEGV} genuino de la
aplicación y reenviar los ajenos al manejador que correspondiera, porque
tragárselos convertiría un error del programa en un cuelgue. El manejador
corre en contexto de señal, donde casi nada es seguro de llamar: no se puede
abrir un socket, ni esperar una respuesta, ni tomar un cerrojo que el hilo
interrumpido pudiera tener. La comunicación tendría que delegarse a un hilo
aparte, con el manejador limitado a despertarlo y esperar, usando únicamente
operaciones permitidas en ese contexto.

Después quedaría el problema que el permiso de página no resuelve: con
\code{mprotect} se detecta el primer acceso de escritura, pero no se sabe
\emph{qué} se escribió sin comparar la página contra una copia previa, ni
cuándo la región crítica terminó, porque ya nadie llama a \code{release}. Eso
obliga a elegir un modelo de consistencia relajado y un punto de
sincronización explícito de todas formas —una barrera, un cerrojo— con lo
cual parte de la transparencia que se buscaba vuelve a la interfaz por otra
puerta.
